% !PRINTED 438
\weberpage{438}
Aus diesen Formeln ergiebt sich leicht nach (5), (8), (14):
\weq{18}{\begin{aligned}
\omega\Theta &= \begin{pmatrix} \gamma, & \alpha\varepsilon^{-2}, & \beta\varepsilon \\ \alpha\varepsilon^2, & \beta, & \gamma\varepsilon^3 \\ \beta\varepsilon^{-1}, & \gamma\varepsilon^{-3}, & \alpha \end{pmatrix}, \quad
\omega\Theta^2 = \begin{pmatrix} \gamma, & \alpha\varepsilon, & \beta\varepsilon^3 \\ \alpha\varepsilon^{-1}, & \beta, & \gamma\varepsilon^2 \\ \beta\varepsilon^{-3}, & \gamma\varepsilon^{-2}, & \alpha \end{pmatrix}, \\
\Theta\omega &= \omega\Theta^3 = \begin{pmatrix} \alpha, & \beta\varepsilon^{-1}, & \gamma\varepsilon^{-3} \\ \beta\varepsilon, & \gamma, & \alpha\varepsilon^{-2} \\ \gamma\varepsilon^3, & \alpha\varepsilon^2, & \beta \end{pmatrix}.
\end{aligned}}

Hiernach ist es nun sehr einfach, die charakteristischen Relationen des Theorems I., \S.\,72 für unsere Gruppe durch wirkliche Ausrechnung zu bestätigen, und damit ist also die ternäre Gruppe 168\textsuperscript{sten} Grades hergestellt\footnote{Vergl. F. \gesperrt{Klein}, \glqq Ueber die Transformation siebenter Ordnung der elliptischen Functionen\grqq{}. \glqq Ueber die Auflösung gewisser Gleichungen vom 7\textsuperscript{ten} und 8\textsuperscript{ten} Grade\grqq{}. Mathem. Annalen, Bd. 14 u. 15. \gesperrt{Klein-Fricke}, Vorlesungen über Modul-Functionen, Bd. I, dritter Abschnitt, Capitel VII. \gesperrt{Gordan}, \glqq Ueber Gleichungen 7\textsuperscript{ten} Grades mit einer Gruppe von 168 Substitutionen\grqq{}. Mathem. Annalen, Bd. 20, 25.}).

\begin{center}
\S.\,116.\\[2pt]
\gesperrt{Pole und Axen der ternären Gruppen.}
\end{center}

Wir wollen uns jetzt der Kürze halber einer geometrischen Ausdrucksweise bedienen, indem wir die Variablen $x_1, x_2, x_3$ als Dreieckscoordinaten eines Punktes $x$ in einer Ebene betrachten, obwohl auch imaginäre Werthe von $x_1, x_2, x_3$ zulässig sein sollen.

Bedeutet $A$ eine ternäre lineare Substitution mit der Determinante 1, und ist
\weq{1}{(y) = A(x),}
so sind $y_1, y_2, y_3$ die Coordinaten eines zweiten Punktes $(y)$, bezogen auf dasselbe Coordinatensystem, der aus $(x)$ durch die Substitution $A$ abgeleitet ist.

Wenn $x$ eine gerade Linie durchläuft, so durchläuft auch $y$ eine gerade Linie, und wenn die Gleichungen dieser beiden Linien
\weq{2}{u_1x_1 + u_2x_2 + u_3x_3 = 0, \quad v_1y_1 + v_2y_2 + v_3y_3 = 0}
sind, so ergiebt sich durch Einsetzen von (1) in (2), dass $u_1, u_2, u_3$ mit $v_1, v_2, v_3$ durch die \gesperrt{transponirte Substitution} von $A$:
\[ (u) = A_1(v) \]
%CONT>>
